\documentclass[10pt,twocolumn,letterpaper]{article}

\providecommand{\texorpdfstring}[2]{#1}
\usepackage{cvpr}
\usepackage{amsmath}
\usepackage{booktabs}
\usepackage{graphicx}
\usepackage{microtype}
\usepackage{xcolor}
\definecolor{cvprblue}{rgb}{0.21,0.49,0.74}
\usepackage[breaklinks,colorlinks,allcolors=cvprblue]{hyperref}

\title{Mid-Term Report: Sensor-Adaptive Agricultural Vision Foundation Models}

\author{
Md Min-Ha-Zul Abedin \and
\href{https://auburn.instructure.com/courses/1739241/users/4120097}{Md Mehedi Hasan}
}

\begin{document}
\maketitle

\begin{abstract}
We study continual self-supervised pretraining of official Vision Transformer
(ViT) checkpoints on a heterogeneous agricultural image corpus. The project
targets transfer across crops, diseases, acquisition conditions, and sensing
modalities while minimizing changes to the pretrained backbone. Since the
proposal, we have implemented a unified framework supporting ImageNet-, MAE-,
DINOv2-, and DINOv3-initialized ViT-S/B/L models; masked image modeling (MIM);
DINO-style student--teacher training; and a learnable $1\!\times\!1$ spectral
adapter for RGB and multiband inputs. The data layer now streams folders, ZIP
and nested-ZIP archives, GeoTIFF, NumPy arrays, and WebDataset shards. A
validated catalog contains 5,175,016 RGB images, substantially exceeding the
412,454-image corpus available at proposal time. Engineering evaluation shows
successful end-to-end MIM and DINO runs, five-band input support, and a
474-epoch SparK--YOLO run whose best validation loss decreased by 18.46\%.
These results establish implementation readiness, but they are not yet evidence
of improved representation quality. Real multispectral coverage, duplicate
control, source-balanced sampling, and downstream held-out-domain comparisons
remain in progress. The next phase will establish ImageNet baselines and measure
linear-probe, few-shot, and fine-tuning performance before scaling pretraining.
\end{abstract}

\section{Introduction}
Agricultural vision systems encounter variation that is less prominent in many
standard image benchmarks: the same crop changes with growth stage, stress,
season, geography, management practice, camera, and illumination. Disease
datasets are often composed of centered leaf photographs, whereas deployment
may involve cluttered field scenes, unmanned aerial vehicle imagery, or
multispectral rasters. Labels for these domains are costly and fragmented among
species identification, disease classification, weed recognition, detection,
segmentation, counting, and temporal monitoring. Consequently, a model that
performs well on a random split of one dataset may transfer poorly to a new
farm or sensor.

Our project asks whether continual self-supervised pretraining can adapt an
official pretrained ViT to agricultural visual structure without discarding the
benefits of generic pretraining. We compare reconstruction-based MIM with
view-invariant DINO-style self-distillation and will test a sequential
MIM$\rightarrow$DINO schedule. A lightweight spectral adapter maps an arbitrary
number of aligned bands to the three-channel interface of an official RGB ViT.
This design keeps the backbone unchanged and makes sensor adaptation directly
ablatable.

At mid-term, the main achievement is a research-ready implementation and a much
larger indexed corpus. The software now supports both objectives, multiple
official checkpoint families, archive-scale streaming, reproducible state, and
downstream task runners. However, the scientific hypothesis remains open:
smoke-test losses and passing unit tests demonstrate correctness, not superior
transfer. We therefore report engineering results separately from scientific
results and reserve performance claims for controlled downstream experiments.

\section{Background and Related Work}
\textbf{Vision Transformers.} ViT represents an image as a sequence of patch
tokens and applies a Transformer encoder, demonstrating that large-scale image
pretraining can rival convolutional networks~\cite{dosovitskiy2021vit}. Our
system uses official ViT-S, ViT-B, and ViT-L implementations rather than a
project-specific transformer, allowing initialization and scale to be controlled
across experiments.

\textbf{Masked image modeling.} MAE masks a large fraction of image patches and
learns to reconstruct their pixels, producing transferable representations with
an asymmetric encoder--decoder design~\cite{he2022mae}. Related masked-image
methods show that reconstruction or token prediction can provide an effective
self-supervised signal~\cite{bao2022beit}. SparK extends this principle to
convolutional encoders by operating on a sparse hierarchy~\cite{tian2023spark}. Agricultural lesions, venation, leaf
edges, and canopy texture motivate testing whether this local reconstruction
signal benefits domain adaptation.

\textbf{Self-distillation.} DINO trains a student to match the centered and
sharpened predictions of an exponential-moving-average (EMA) teacher under
different augmented views~\cite{caron2021dino}. DINOv2 demonstrates that careful
data curation and self-distillation can produce robust visual features without
labels~\cite{oquab2024dinov2}. Reconstruction and self-distillation emphasize
different invariances, motivating MIM-only, DINO-only, and sequential training
under a shared data and backbone stack.

\textbf{Agricultural transfer and spectral adaptation.} PlantVillage helped
establish deep learning for image-based plant-disease recognition, but its
controlled backgrounds also illustrate the difference between benchmark and
field conditions~\cite{hughes2015plantvillage,mohanty2016using}. Remote-sensing
foundation models such as SatMAE adapt masked pretraining to temporal and
multispectral satellite imagery~\cite{cong2022satmae}. In contrast, our adapter
preserves an official RGB backbone and learns a small sensor-specific projection,
which enables a controlled comparison against RGB-only input and direct
patch-embedding expansion.

\section{Approach}
\subsection{Data and Sampling}
The proposal reported 412,454 images from 13 downloaded datasets. The current
validated catalog indexes 5,175,016 loader-visible images across approximately
1~TB of ZIP and TAR archives. Sources cover plant species, diseases, seedlings,
weeds, and crop imagery, but the present catalog remains overwhelmingly RGB.
Each catalog record stores a stable sample key, source identity, and portable
archive-member path. Large sources are converted to roughly 1~GB WebDataset
shards after RGB conversion, bounded resizing, and validation. Parallel workers
process one archive at a time to reduce archive-open and seek overhead.

To prevent a large species collection from dominating training, the planned
sampler first selects a source and then samples within that source. Validation
splits are group-disjoint when source metadata permits. Before final experiments,
we will freeze the manifest, record licenses and versions, and remove exact and
near-duplicate clusters before creating train and held-out-domain partitions.

\begin{table}[t]
\centering
\caption{Corpus status at proposal and mid-term. Counts describe indexed files,
not unique biological observations.}
\label{tab:corpus}
\small
\begin{tabular}{lrr}
\toprule
Property & Proposal & Mid-term \\
\midrule
Validated images & 412,454 & 5,175,016 \\
Local storage & 35.86 GB & $\sim$1 TB \\
Primary modality & RGB & RGB \\
Real multispectral archive & No & No \\
Archive streaming & Audit only & ZIP/TAR/WDS \\
\bottomrule
\end{tabular}
\end{table}

\subsection{Sensor-Adaptive Backbone}
For an input $\mathbf{x}\in\mathbb{R}^{C\times H\times W}$, a learnable
$1\!\times\!1$ convolution projects the $C$ sensor bands into the three channels
expected by an official pretrained ViT:
\begin{equation}
 \mathbf{z}_{h,w}=A\mathbf{x}_{h,w}+\mathbf{b},\qquad
 A\in\mathbb{R}^{3\times C}.
 \label{eq:adapter}
\end{equation}
For RGB, $A$ is initialized to the $3\times3$ identity and $\mathbf{b}=0$.
For multiband input, available RGB-like channels initialize the corresponding
outputs; weights for additional bands begin at zero but remain trainable. The
adapted image is normalized with the statistics associated with the selected
official checkpoint, then patchified by the unchanged ViT. The adapter adds only
$3C+3$ parameters and cleanly separates sensor mixing from spatial feature
learning.

\subsection{Self-Supervised Objectives}
In MIM, a random set $\mathcal{M}$ masks 75\% of patch tokens. A linear
reconstruction head predicts the adapted RGB patch $\mathbf{p}_i$ from encoded
token $\mathbf{h}_i$; loss is evaluated only at masked positions:
\begin{equation}
 \mathcal{L}_{\mathrm{MIM}}=
 \frac{1}{|\mathcal{M}|}\sum_{i\in\mathcal{M}}
 \left\|g(\mathbf{h}_i)-\mathbf{p}_i\right\|_2^2.
 \label{eq:mim}
\end{equation}

For DINO-style training, the student receives global and local crops while the
teacher receives paired global crops. With student temperature $\tau_s$, teacher
temperature $\tau_t$, and running center $\mathbf{c}$, the cross-view loss is
\begin{equation}
 \mathcal{L}_{\mathrm{DINO}}=-\sum_k
 q_t^{(k)}\log q_s^{(k)},\quad
 q_t=\mathrm{softmax}((\mathbf{o}_t-\mathbf{c})/\tau_t).
 \label{eq:dino}
\end{equation}
The teacher parameters, including adapter, backbone, and projection head, are
updated by $\theta_t\leftarrow m\theta_t+(1-m)\theta_s$. The framework supports
ImageNet, MAE, DINOv2, and DINOv3 initialization and an explicit
MIM$\rightarrow$DINO handoff that starts a new optimization run from a compatible
MIM checkpoint.

\subsection{Training and Evaluation Protocol}
Training uses BF16 mixed precision, gradient accumulation, optional gradient
checkpointing, cosine scheduling, and multi-process data loading. Every run
stores resolved configuration, environment and Git metadata, metrics, figures,
best/latest checkpoints, optimizer state, and random-number-generator state.
WebDataset provides CPU streaming; NVIDIA DALI is available on Linux for GPU
JPEG decoding. Development began on a single RTX~4090, and the next stage will
use Auburn's Slurm cluster for longer controlled runs.

Scientific evaluation will compare the unadapted official checkpoint, RGB-only
continual pretraining, MIM, DINO, sequential MIM$\rightarrow$DINO, frozen versus
trainable adapters, and direct patch-embedding expansion. Classification will
report accuracy and macro-F1; detection, segmentation, and counting will use
mAP, IoU/Dice, and MAE/RMSE as appropriate. Linear probing, few-shot subsets,
and full fine-tuning will be repeated across at least three seeds. The primary
test sets will hold out whole sources, crops, geographies, seasons, or sensors,
with mean, standard deviation, and paired comparisons reported where splits are
shared.

\section{Intermediate and Preliminary Results}
\subsection{Verified Engineering Results}
Table~\ref{tab:results} summarizes completed checks. End-to-end runs produced
checkpoints, metrics, training curves, MIM reconstructions, DINO augmented views,
and student--teacher similarity matrices. Five-band float32 arrays passed through
the learnable adapter and official ImageNet ViT-S. DINO also ran from an official
DINOv2 ViT-S checkpoint with its patch-14 geometry, and a compatible MIM
checkpoint successfully initialized a new five-band DINO run. Two-epoch resume
tests restored optimizer, scheduler, scaler, history, RNG, and loader state.

\begin{table}[t]
\centering
\caption{Mid-term implementation validation. ``Pass'' indicates successful
execution and expected artifacts, not a downstream accuracy result.}
\label{tab:results}
\small
\begin{tabular}{p{0.48\linewidth}p{0.18\linewidth}p{0.19\linewidth}}
\toprule
Check & Input & Result \\
\midrule
MIM, official ImageNet ViT-S & 5-band NPY & Pass \\
DINO, official ImageNet ViT-S & RGB PNG & Pass \\
DINO, official DINOv2 ViT-S & RGB PNG & Pass \\
MIM$\rightarrow$DINO initialization & 5-band NPY & Pass \\
Two-epoch checkpoint resume & RGB/5-band & Pass \\
Automated test suite & 124 tests & 124 pass \\
Shard builder smoke test & ZIP/plain & 2 pass \\
\bottomrule
\end{tabular}
\end{table}


\subsection{Long-Run SparK--YOLO Pretraining}
We also evaluated the convolutional branch in a 500-epoch SparK run using the
WebDataset pipeline and a YOLO11-L backbone. Initialization loaded the official
YOLO checkpoint, inserted $1\!\times\!1$ expansion bridges where dense and sparse
channel dimensions differed, and transferred 215 tensors into the sparse
encoder. Each epoch comprised 4,000 training batches. The learning rate warmed
from $2\!\times\!10^{-6}$ to $10^{-5}$ over the first five epochs, remained near
that value through the early run, and then followed a cosine decay toward
$10^{-6}$.

At the time of analysis, 474 epochs were complete and epoch 475 had processed
1,232 of 4,000 batches. Validation loss fell from 0.8071 at epoch 1 to a best of
0.6581 at epoch 472, an 18.46\% relative reduction. The corresponding training
loss was 0.6524, leaving a gap of only 0.0057. Table~\ref{tab:spark-run} shows
representative checkpoints along the trajectory. Most improvement occurred in
the first 100 epochs; progress thereafter was slower but persistent, with new
minima at epochs 168, 257, 297, 353, 411, 427, 452, and 472. Over the final 25
complete epochs, validation loss averaged 0.6610 with a standard deviation of
0.0014. This narrow band indicates a stable late-stage plateau rather than
divergence. The last observed validation loss, 0.6615, was slightly above the
minimum, so downstream experiments will use the saved best checkpoint rather
than the terminal state.

\begin{table}[t]
\centering
\caption{Selected points from the SparK--YOLO pretraining log. Epoch 475 was
still running and is therefore excluded.}
\label{tab:spark-run}
\small
\begin{tabular}{rrrr}
\toprule
Epoch & Train & Validation & LR \\
\midrule
1   & 0.8250 & 0.8071 & $2\!\times\!10^{-6}$ \\
100 & 0.6782 & 0.6808 & $9\!\times\!10^{-6}$ \\
297 & 0.6601 & 0.6598 & $4\!\times\!10^{-6}$ \\
472 & 0.6524 & \textbf{0.6581} & $1\!\times\!10^{-6}$ \\
474 & 0.6431 & 0.6615 & $1\!\times\!10^{-6}$ \\
\bottomrule
\end{tabular}
\end{table}
The automated suite covers RGB, NPY, GeoTIFF, ZIP and nested-ZIP loading;
group-disjoint splits; NoData handling; undersized-image padding; official model
selection; adapter initialization; MIM and DINO tensor behavior; EMA updates;
paired crop replay; strict configuration parsing; atomic checkpointing; and
RoPE-based DINOv3 backbones. Recent shard-builder tests additionally verify
parallel archive grouping and valid TAR output. These checks reduce the risk
that a long pretraining run fails because of data format, checkpoint, or model
geometry errors.

\subsection{Interpretation and Remaining Gaps}
The corpus expansion, successful pipelines, and four-day SparK--YOLO run satisfy
the mid-term engineering milestone: the planned models can now train
reproducibly at archive scale. The long-run loss trajectory is meaningful
evidence that sparse conversion, transferred YOLO weights, checkpoint selection,
and WebDataset streaming remain numerically stable beyond smoke-test duration.
It is not evidence that the learned representation improves a recognition task:
reconstruction loss measures the pretext objective and can decline without
improving transfer.
The results also support the feasibility of the spectral interface because the
same official backbone accepts both RGB and five-band tensors without replacing
its patch embedding. They do \emph{not}, however, test whether additional bands
improve predictions; the five-band data used for verification are synthetic,
and the catalog lacks a verified real multispectral archive.

No model is yet shown to outperform its official ImageNet or DINO initialization
on a held-out agricultural benchmark. We therefore make no claim of improved
accuracy, robustness, or label efficiency at this stage. The largest immediate
risks are duplicate leakage, imbalance from the 306k-image PlantNet subset and
other large sources, RGB-only bias, and I/O saturation during multi-GPU training.
The first three are addressed by manifest freezing, duplicate clustering,
source-balanced sampling, and external domain splits; the last will be measured
by comparing CPU WebDataset and DALI throughput before requesting more GPUs.

The next milestone is a ViT-S experiment matrix on frozen data: official
initialization, MIM, DINO, and sequential MIM$\rightarrow$DINO, followed by
linear probing and low-label fine-tuning on source-disjoint disease and species
benchmarks. Only recipes that improve repeated-seed transfer will scale to ViT-B.
In parallel, we will add a licensed real multispectral source and evaluate
frozen/trainable adapters against direct patch-embedding expansion. This order
prioritizes evidence over model size and keeps the final conclusions aligned with
the data actually available.

{\small
\begin{thebibliography}{9}
\bibitem{dosovitskiy2021vit}
A. Dosovitskiy et al. An Image Is Worth 16x16 Words: Transformers for Image
Recognition at Scale. \textit{ICLR}, 2021.

\bibitem{he2022mae}
K. He et al. Masked Autoencoders Are Scalable Vision Learners.
\textit{CVPR}, 2022.

\bibitem{bao2022beit}
H. Bao et al. BEiT: BERT Pre-Training of Image Transformers.
\textit{ICLR}, 2022.

\bibitem{tian2023spark}
K. Tian, Y. Jiang, Q. Diao, C. Lin, L. Wang, and Z. Yuan. Designing BERT for
Convolutional Networks: Sparse and Hierarchical Masked Modeling. \textit{ICLR}, 2023.

\bibitem{caron2021dino}
M. Caron et al. Emerging Properties in Self-Supervised Vision Transformers.
\textit{ICCV}, 2021.

\bibitem{oquab2024dinov2}
M. Oquab et al. DINOv2: Learning Robust Visual Features without Supervision.
\textit{TMLR}, 2024.

\bibitem{hughes2015plantvillage}
D. Hughes and M. Salath\'e. An Open Access Repository of Images on Plant Health
to Enable the Development of Mobile Disease Diagnostics. \textit{CoRR}, 2015.

\bibitem{mohanty2016using}
S. Mohanty, D. Hughes, and M. Salath\'e. Using Deep Learning for Image-Based
Plant Disease Detection. \textit{Frontiers in Plant Science}, 2016.

\bibitem{cong2022satmae}
Y. Cong et al. SatMAE: Pre-training Transformers for Temporal and Multi-Spectral
Satellite Imagery. \textit{NeurIPS}, 2022.
\end{thebibliography}
}

\end{document}
